\documentclass[11pt,a4paper]{article}
\usepackage[utf8]{inputenc}
\usepackage[T1]{fontenc}
\usepackage{lmodern}
\usepackage{geometry}
\usepackage{amsmath,amssymb}
\usepackage{hyperref}
\usepackage{enumitem}
\geometry{margin=1in}

\title{TRM V3.1--V3.2 Addendum:\\From Memory Action-Closure to Lattice-Derived Minimal Action}
\author{
Fabrice Wieser\\
Independent Researcher\\
\url{https://github.com/MagusDraconis/TRM_Cosmology}
}
\date{
Technical Addendum to the TRM/TQM V3.0 Review Baseline\\
June 29, 2026
}

\begin{document}
\maketitle

\noindent
This document is a compact, citable supplement to the V3.0 trilogy, not a new major paper generation.

\section{Abstract}
This addendum provides a claim-safe update to the TRM V3.0 paper trilogy.
V3.1 strengthens Gap 1 to an \emph{action-derived memory-closure candidate}, and V3.2 strengthens the action layer to a \emph{minimal effective action that is plausibly lattice-derived under comparative guard tests}.
These results improve structural and comparative support while remaining below theorem-level first-principles closure.

\section{Relation to V3.0 Review Baseline}
V3.0 established the public review baseline with broad weak-field tested-effective guard coverage.
It did not claim theorem-level microscopic closure and did not claim GR replacement.
V3.1 and V3.2 are targeted hardening layers on top of that baseline, not a replacement of the baseline claim boundary and not a ``V4'' reset.

\section{V3.1: Action-Derived Memory-Closure Candidate}
V3.1 focuses on Gap 1 and hardens the memory-channel path:
\[
A_{\mathrm{dyn}} \propto \phi
\rightarrow
A_{\mathrm{dyn}}^2 \lvert \dot{\mu} \rvert
\rightarrow
\phi^2 \lvert \dot{\mu} \rvert.
\]

\noindent
Status in V3.1:
\begin{itemize}[leftmargin=1.5em]
\item Gap 1 is strengthened to an \emph{action-derived memory-closure candidate}.
\item The path is guard-supported and variation-compatible.
\item It remains \emph{not theorem-level first-principles closure}.
\end{itemize}

\section{V3.2: Minimal Effective Action from TQM Lattice/Synchronization Principles}
V3.2 moves one level deeper by testing whether the minimal effective action itself is plausibly motivated by TQM lattice/synchronization structure.

\noindent
Status in V3.2:
\begin{itemize}[leftmargin=1.5em]
\item The minimal effective action is \emph{plausibly lattice-derived under comparative guard tests}.
\item This remains a candidate-level derivation status, not theorem-level closure.
\end{itemize}

\section{Guard Summary}
V3.2 guard block (UA16--UA23), in compact form:
\begin{itemize}[leftmargin=1.5em]
\item \textbf{UA16}: lattice energy reduces to minimal scalar action in tested weak-field cases.
\item \textbf{UA17}: $A_{\mathrm{dyn}}^2\kappa$ is preserved under coarse-graining.
\item \textbf{UA18}: V3.1 UF13--UF15 behavior is reproduced without retuning.
\item \textbf{UA19}: non-minimal terms remain penalized/subleading.
\item \textbf{UA20}: MC/FD/TO limits remain preserved.
\item \textbf{UA21}: alternative lattice-energy forms degrade minimal-action closure quality.
\item \textbf{UA22}: minimal-action closure remains robust under controlled lattice perturbations.
\item \textbf{UA23}: non-minimal lattice actions do not outperform the minimal action under penalty-aware comparison.
\end{itemize}

\noindent
Combined interpretation: the framework supports not only positive-fit behavior of one action form, but comparative preference and robustness for the minimal action class in the tested regime.

\section{Remaining Theorem-Level Gaps}
The following remain open:
\begin{enumerate}[leftmargin=1.5em]
\item theorem-level first-principles microscopic derivation of the minimal effective action,
\item theorem-level uniqueness/necessity of the selected leading interaction hierarchy,
\item full theorem-level closure across remaining first-principles gaps.
\end{enumerate}

\section{Conclusion}
Relative to V3.0, V3.1 and V3.2 provide materially stronger claim-safe support for the memory/action layer:
\begin{itemize}[leftmargin=1.5em]
\item V3.1: action-derived memory-closure candidate,
\item V3.2: minimal effective action plausibly lattice-derived under comparative guards.
\end{itemize}

\noindent
Boundary remains explicit:
\begin{itemize}[leftmargin=1.5em]
\item no GR replacement claim,
\item no theorem-level first-principles closure claim.
\end{itemize}

\end{document}
